\chapter{Parent-origin бинарного \texorpdfstring{$u/d$}{u/d}-селектора incidence}
\label{ch:v8-extra-mass-up-down-binary-selector-origin}

\section{Точный бинарный остаток}

Предыдущий гейт оставил два назначения
\begin{equation}
 \mathcal M_{ud}=\{B\mapsto u,\ M\mapsto Y;\quad
                    B\mapsto d,\ M\mapsto Y\}.
 \label{eq:v8-up-down-selector-binary-menu}
\end{equation}
Они имеют одинаковые размерности и одинаковые не абелевы индексы, но
разные гиперзаряды:
\begin{equation}
 d_u=d_d=3,\qquad T_2(u)=T_2(d)=0,\qquad
 T_3(u)=T_3(d)=\frac12,\qquad
 y_u=\frac53,\quad y_d=\frac23.
 \label{eq:v8-up-down-selector-representation-data}
\end{equation}
Следовательно, размерность, цвет и слабый тип бинарную свободу не снимают.

\section{Квадрат гиперзаряда как условный различитель}

Комплексные квадратичные $U(1)$-индексы равны
\begin{equation}
 I_1(u)=3y_u^2=\frac{25}{3},\qquad
 I_1(d)=3y_d^2=\frac{4}{3},\qquad
 I_1(u)-I_1(d)=7.
 \label{eq:v8-up-down-selector-u1-index-gap}
\end{equation}
Поэтому условный дискретный функционал
\begin{equation}
 S_\alpha(e)=\alpha I_1(e),\qquad e\in\{u,d\},
 \label{eq:v8-up-down-selector-conditional-score}
\end{equation}
даёт
\begin{equation}
 \arg\min S_\alpha=
 \begin{cases}
 \{d\},&\alpha>0,\\
 \{u,d\},&\alpha=0,\\
 \{u\},&\alpha<0.
 \end{cases}
 \label{eq:v8-up-down-selector-score-minimizers}
\end{equation}
Таким образом, положительный $Y^2$-вес условно предпочитает $d$-ветвь.
Но это ещё не parent-origin: необходимо вывести сам typed-вес, его знак и
правило вариации по диаграммам.

\section{Real-след удаляет нечётный кандидат}

На частице и Real-сопряжении введём
\begin{equation}
 Y_{\mathbb R}=\operatorname{diag}
 \left(\frac53I_3,-\frac53I_3,
       \frac23I_3,-\frac23I_3\right),
 \quad
 A_{ud}=\operatorname{diag}(I_3,I_3,-I_3,-I_3).
 \label{eq:v8-up-down-selector-real-hypercharge}
\end{equation}
Все нечётные Real-моменты, которые могли бы нести ориентированный знак,
сокращаются попарно:
\begin{equation}
 \Tr(A_{ud}Y_{\mathbb R})=0,
 \qquad
 \Tr(A_{ud}Y_{\mathbb R}^{3})=0.
 \label{eq:v8-up-down-selector-odd-real-cancellation}
\end{equation}
Это тот же механизм, который делает каждую добавленную вектороподобную пару
аномально нейтральной. Чётный момент различает типы:
\begin{equation}
 \Tr(A_{ud}Y_{\mathbb R}^{2})
 =6\left(\frac{25}{9}-\frac49\right)=14.
 \label{eq:v8-up-down-selector-even-real-discriminator}
\end{equation}
Однако обычное спектральное действие является функционалом уже выбранной
конечной геометрии; оно не превращает две диаграммы Краевского в две точки
одного полевого многообразия \cite{v8-up-down-selector-krajewski,
v8-up-down-selector-spectral-action}.

\section{Почему сравнение двух графов не есть уравнение движения}

Пусть $P_u,P_d$ --- Real-проекторы двух возможных цветных incidence-концов.
Они ортогональны, имеют одинаковый ранг и конечное расстояние:
\begin{equation}
 P_uP_d=0,\qquad \Tr P_u=\Tr P_d=6,\qquad
 \Tr(P_u-P_d)^2=12.
 \label{eq:v8-up-down-selector-projector-distance}
\end{equation}
В текущем пространстве полей эти проекторы являются структурными данными,
а не динамическими переменными:
\begin{equation}
 \dim T_{P_u}\mathcal M_{ud}
 =\dim T_{P_d}\mathcal M_{ud}=0,
 \qquad \frac{\delta S}{\delta P_e}\ \text{не определено}.
 \label{eq:v8-up-down-selector-zero-tangent}
\end{equation}
Поэтому разность $S_\alpha(u)-S_\alpha(d)=7\alpha$ сравнивает две разные
модели. Для физического выбора потребовались бы сумма по геометриям, мера
или единый больший оператор, внутри которого обе incidence-компоненты
являются состояниями одной динамической переменной.

\section{Минимальный недостающий typed-оператор}

Квадрат гиперзаряда канонически задаёт центрированный различитель
\begin{equation}
 S_{ud}=Y_{\mathbb R}^{2}-\frac{29}{18}I
 =\frac76P_u-\frac76P_d.
 \label{eq:v8-up-down-selector-centered-discriminator}
\end{equation}
Его нормировка и спаривание с бинарным направлением фиксированы:
\begin{equation}
 \Tr S_{ud}^{2}=\frac{49}{3},
 \qquad \Tr(A_{ud}S_{ud})=14.
 \label{eq:v8-up-down-selector-discriminator-norm}
\end{equation}
Но текущий parent не содержит ни динамического бинарного носителя, ни
связи этого носителя с $S_{ud}$. Более того, знак такой связи меняет
предпочитаемую ветвь:
\begin{equation}
 V_{\rm sel}(P_e)=g_{ud}\Tr(S_{ud}P_e),\qquad
 g_{ud}\mapsto-g_{ud}\ \Longleftrightarrow\ u\leftrightarrow d.
 \label{eq:v8-up-down-selector-missing-coupling-sign}
\end{equation}
Именно происхождение typed-связи и её знака, а не вычисление $Y^2$, является
следующим обязательством.

Итоговые реестры:
\begin{equation}
 N_{\rm exact\ representation}=8/8,\qquad
 N_{\rm conditional\ discriminator}=5/5,\qquad
 N_{\rm parent\ origin}=0/7.
 \label{eq:v8-up-down-selector-ledgers}
\end{equation}

\begin{theorem}[Условный гиперзарядный различитель без parent-origin]
Квадрат гиперзаряда точно различает два финальных incidence-назначения:
их комплексные индексы равны $25/3$ и $4/3$. Положительный условный вес
предпочитает $d$-ветвь. Однако нечётные Real-моменты сокращаются, incidence-
проекторы не являются переменными текущего действия, а typed-связь с
$Y^2$ и её знак не наследованы. Поэтому бинарный выбор не выведен.
\end{theorem}

\begin{nogocriterion}[Запрет минимизации между фиксированными диаграммами]
Нельзя объявлять $B\mapsto d$ физически выбранным только потому, что
$I_1(d)<I_1(u)$. Такое сравнение становится вариационным принципом лишь
после задания общего пространства геометрий и его меры либо единого
динамического носителя бинарного incidence. Нельзя также импортировать знак
положительного gauge-кинетического коэффициента в отсутствующий портальный
член endpoint-действия.
\end{nogocriterion}

\begin{protocol}\begin{enumerate}
\item финальных incidence-кандидатов: \textbf{два};
\item размерность, $SU(2)$ и $SU(3)$ различают их: \textbf{нет};
\item $U(1)$-индексы: \textbf{$25/3$ и $4/3$};
\item положительный $Y^2$-score условно выбирает: \textbf{$d$};
\item нечётные Real-моменты дают селектор: \textbf{нет};
\item вектороподобная аномалия даёт селектор: \textbf{нет};
\item incidence-проектор динамичен: \textbf{нет};
\item межгеометрическая мера наследована: \textbf{нет};
\item центрированный typed-различитель: \textbf{$S_{ud}$ построен};
\item знак его связи выведен: \textbf{нет};
\item exact representation: \textbf{$8/8$};
\item conditional discriminator: \textbf{$5/5$};
\item parent-origin: \textbf{$0/7$};
\item следующий гейт: \textbf{минимальная typed-архитектура различителя}.
\end{enumerate}\end{protocol}

Машинный сертификат сохранён в
\path{s2t_v8_baryon_c0_singlet_triplet_central_gap_isotypic_channel_extra_edge_mass_two_source_cubic_incidence_up_down_binary_selector_parent_origin_gate_results.json}.

\begin{thebibliography}{9}
\bibitem{v8-up-down-selector-krajewski} T.~Krajewski,
\emph{Classification of Finite Spectral Triples}, arXiv:hep-th/9701081.
\bibitem{v8-up-down-selector-spectral-action} A.~H.~Chamseddine,
A.~Connes, \emph{The Spectral Action Principle}, arXiv:hep-th/9606001.
\end{thebibliography}